\section{Explicit analytic identities with full constants}
\label{sec:identities}

The formal classification establishes existence of reductions but does
not, by itself, print their logarithmic corrections. This section does
so for an infinite family and for the remaining exceptional ratio
$3/5$. All the identities in this section are ordinary analytic
equalities. No integer-relation search is used in their proofs.

\subsection{Classical normalized functional equations}

For $0<t<1$ define the Rogers dilogarithm
\begin{equation}\label{eq:Rogers}
 \Rog(t)=\Li_2(t)+\frac12\log t\log(1-t),
 \qquad \Rog(1/2)=\frac{\pi^2}{12}.
\end{equation}
We use its reflection and five-term identities in the real interval:
\begin{align}
 \Rog(t)+\Rog(1-t)&=\frac{\pi^2}{6},\label{eq:Rog-reflect}\\
 \Rog(A)+\Rog(B)&=\Rog(AB)
  +\Rog\left(\frac{A(1-B)}{1-AB}\right)
  +\Rog\left(\frac{B(1-A)}{1-AB}\right),
 \quad 0<A,B<1.\label{eq:Rog-five}
\end{align}
These are classical; their normalization and relation to the ordinary
dilogarithm are reviewed in \cite{ZagierDilog,RZ}.

Put
\begin{equation}\label{eq:modified-F}
 f(x)=F(x)-F(1)+\frac{\pi^2}{12}(x-2+x^{-1})-\frac14\log^2x.
\end{equation}
The functional equations needed below are the normalized inversion and
three-term equations and the determinant-$2$ Hecke equation of
\cite[(15), (22a)]{RZ}:
\begin{align}
 f(1/x)&=-f(x),\label{eq:f-inversion}\\
 f(x)-f(x+1)-f\left(\frac{x}{x+1}\right)
 &=\Rog(1/2)-\Rog\left(\frac{x}{x+1}\right),\label{eq:f-three}\\
 \begin{split}
 f(2x)+f(x/2)+f\left(\frac{x+1}{2}\right)
 +f\left(\frac{2x}{x+1}\right)-3f(x)
 ={}&\Rog\left(\frac{x}{x+1}\right)-\Rog(1/2)\\
 &+\tfrac12\log2\log x.
 \end{split}\label{eq:f-Hecke}
\end{align}
They hold for $x>0$ with real logarithms. Formula \eqref{eq:J-F}
becomes
\begin{equation}\label{eq:J-f}
 J(x)=f(2x)-2f(x)+f(x/2)+\frac12\log^22
       -\frac{\pi^2}{24}(x-x^{-1}).
\end{equation}

Define the finite log-sine sum
\begin{equation}\label{eq:S}
 S(m)=\sum_{k=1}^{m-1}\log^2\left(2\sin\frac{\pi k}{m}\right),
 \qquad S(1)=0.
\end{equation}
Every sine argument in this sum is positive. The integer formula of
\cite[(23)]{RZ} simplifies under \eqref{eq:modified-F} to
\begin{equation}\label{eq:f-integer}
 f(m)=\frac{(m-1)\pi^2}{24}+\frac12S(m)+\frac14\log^2m,
 \qquad m\ge1.
\end{equation}
These classical inputs are kept separate from the deductions that follow.

\subsection{Every neighboring rational argument}

\begin{theorem}[Uniform neighboring-ratio identity]\label{thm:neighbor}
For every integer $n\ge1$,
\begin{equation}\label{eq:neighbor}
 \boxed{\begin{aligned}
 J\left(\frac n{n+1}\right)={}&\frac12\log^22
 +\frac12\bigl(S(n)-S(n+1)+S(2n+2)-S(2n)\bigr)\\
 &-\frac{\pi^2(n^2-n-1)}{24n(n+1)}.
 \end{aligned}}
\end{equation}
The inverse ratios are evaluated by \eqref{eq:J-reciprocity}.
\end{theorem}
\begin{proof}
Solve \eqref{eq:f-Hecke} for the second difference in
\eqref{eq:J-f}. Set
\[
 x=\frac n{n+1},\quad A=\frac{2n+1}{2n+2},\quad
 B=\frac{2n}{2n+1},\quad C=\frac n{2n+1}.
\]
The three $f$-terms on the resulting right side are $f(x)-f(A)-f(B)$.
For any integer $m\ge1$, \eqref{eq:f-three} gives
\[
 f\left(\frac m{m+1}\right)=f(m)-f(m+1)-\Rog(1/2)
                                      +\Rog\left(\frac m{m+1}\right).
\]
The integer terms telescope, leaving
\begin{equation}\label{eq:neighbor-intermediate}
 \begin{split}
 f(2x)-2f(x)+f(x/2)={}&f(n)-f(n+1)+f(2n+2)-f(2n)\\
 &+\Rog(x)-\Rog(A)-\Rog(B)+\Rog(C)
   +\tfrac12\log2\log x.
 \end{split}
\end{equation}
In \eqref{eq:Rog-five} with these $A,B$, the three right-hand arguments
are respectively $x$, $1/2$ and $C$. Hence the Rogers combination in
\eqref{eq:neighbor-intermediate} is exactly $-\Rog(1/2)$.

Substitution of \eqref{eq:f-integer} contributes $\pi^2/24$ from the
four integer terms. Their ordinary log-square contribution is
\[
 \tfrac14\bigl(\log^2n-\log^2(n+1)+\log^2(2n+2)-\log^2(2n)\bigr)
 =-\tfrac12\log2\log x,
\]
which cancels the cross term. The remaining $\pi^2$ coefficient in
\eqref{eq:J-f} is
\[
 \frac1{24}-\frac1{12}-\frac1{24}(x-x^{-1})
 =-\frac{n^2-n-1}{24n(n+1)}.
\]
This proves \eqref{eq:neighbor}. All Rogers arguments used in the proof
are strictly between zero and one, so no branch adjustment or unspecified
constant is hidden in the cancellation.
\end{proof}

\begin{corollary}[Small explicit logarithmic evaluations]\label{cor:small-J}
Let $\phi=(1+\sqrt5)/2$. Then
\begin{align}
 J(4/3)&=\frac78\log^22-\frac12\log^2(1+\sqrt2)
                     +\frac{5\pi^2}{288},\label{eq:J43}\\
 J(4/5)&=\frac78\log^22+2\log^2\phi
              -\frac12\log^2(1+\sqrt2)-\frac{11\pi^2}{480}.
 \label{eq:J45}
\end{align}
\end{corollary}
\begin{proof}
Write $a=\log2$, $b=\log3$, $c=\log5$, $d=\log\phi$ and
$u=\log(1+\sqrt2)$. Elementary trigonometric products give
\begin{align*}
 S(3)&=\tfrac12b^2,&S(4)&=\tfrac32a^2,&
 S(5)&=\tfrac14c^2+d^2,\\
 S(6)&=a^2+\tfrac12b^2,&S(8)&=\tfrac74a^2+u^2,&
 S(10)&=a^2+\tfrac14c^2+5d^2.
\end{align*}
For instance $2\sin(\pi/10)=\phi^{-1}$ and
$2\sin(3\pi/10)=\phi$, while the two distinct primitive positive
logarithms at conductor $8$ have sum $a/2$ and difference $u$.
Substitute into \eqref{eq:neighbor} for $n=3,4$, using reciprocity for
$n=3$.
\end{proof}

At $n=1,2$, the theorem also recovers
$J(1/2)=\pi^2/48+\log^22/4$ and
$J(2/3)=-\pi^2/144+3\log^22/4$. These are consistency checks against
known special cases, not additional novelty claims.

\subsection{The exceptional value at \texorpdfstring{$3/5$}{3/5}}

\begin{theorem}[A compact exceptional dilogarithm identity]\label{thm:J35}
With $\phi=(1+\sqrt5)/2$,
\begin{equation}\label{eq:J35}
 \boxed{\begin{aligned}
 J(3/5)={}&\Li_2(1/3)-\frac12\Li_2(1/5)
  +\frac12\log^22+\frac12\log^23-\frac14\log^25\\
 &+\log^2\phi-\frac{7\pi^2}{180}.
 \end{aligned}}
\end{equation}
In particular $[1/3]-\tfrac12[1/5]$ is a rational-dilogarithm core.
\end{theorem}
\begin{proof}
We first track the normalized Herglotz terms, then remove the remaining
Rogers terms by a single five-term substitution. Set
$a=\log2$, $b=\log3$, $c=\log5$, $d=\log\phi$.
The classical $F(2/5)$ identity \cite[(24)]{RZ} is
\[
 F(2/5)-F(1)=\tfrac12\Li_2(4/5)-\tfrac15\pi^2
 +\log^2(2\sin(\pi/5))+\log^2(2\sin(2\pi/5))
 +a\log(2/5).
\]
The two log-sine squares sum to $c^2/8+d^2/2$. Converting to $f$ and
to the Rogers normalization gives
\begin{equation}\label{eq:f25}
 f(2/5)=\tfrac12\Rog(4/5)-\tfrac18\pi^2
                      +\tfrac34a^2-\tfrac38c^2+\tfrac12d^2.
\end{equation}
The integer formula gives
\begin{align*}
 f(2)&=\pi^2/24+3a^2/4,&
 f(3)&=\pi^2/12+b^2/2,\\
 f(4)&=\pi^2/8+7a^2/4,&
 f(5)&=\pi^2/6+3c^2/8+d^2/2.
\end{align*}
Therefore
\begin{equation}\label{eq:f25-collapse}
 -f(2)+f(5)+f(2/5)=d^2+\tfrac12\Rog(4/5).
\end{equation}
Inversion and the three-term relation at $3/2$ give
\[
 f(3/5)=-f(2)+f(3)+f(2/5)-\Rog(2/3)+\Rog(3/5).
\]
Also
$f(4/5)=f(4)-f(5)-\Rog(1/2)+\Rog(4/5)$ and
$f(3/4)=f(3)-f(4)-\Rog(1/2)+\Rog(3/4)$.
The Hecke relation at $x=3/5$, together with
\eqref{eq:J-f} and \eqref{eq:f25-collapse}, now becomes
\begin{equation}\label{eq:J35-Rog}
 J(3/5)=d^2+T+\tfrac12a(b-c)+\tfrac12a^2+\frac{2\pi^2}{45},
\end{equation}
where
\[
 T=\Rog(1/2)-\Rog(2/3)+\Rog(3/5)-\tfrac12\Rog(4/5)
                         -\Rog(3/4)+\Rog(3/8).
\]
The substitution $A=1/2$, $B=3/4$ into \eqref{eq:Rog-five} gives
\[
 \Rog(1/2)+\Rog(3/4)=\Rog(3/8)+\Rog(1/5)+\Rog(3/5).
\]
Combining this with reflection at $1/3$ and $1/5$ yields
\[
 T=\Rog(1/3)-\tfrac12\Rog(1/5)-\frac{\pi^2}{12}.
\]
Finally,
\[
 \Rog(1/3)-\tfrac12\Rog(1/5)
 =\Li_2(1/3)-\tfrac12\Li_2(1/5)
  -\tfrac12ab+\tfrac12b^2+\tfrac12ac-\tfrac14c^2.
\]
The cross terms cancel in \eqref{eq:J35-Rog};
$2/45-1/12=-7/180$ gives the stated constant.

To check the core directly, take its boundary:
\[
 \partial\bigl([1/3]-\tfrac12[1/5]\bigr)
 =\cl2\wedge\cl3-\cl2\wedge\cl5
 =\cl2\wedge\cl{3/5},
\]
which agrees with \eqref{eq:odd-boundary}. The analytic argument above,
not this last boundary check alone, supplies the complete constants.
\end{proof}

\subsection{All six non-logarithmic formal exceptions are explicit}

The known identities for the other two reciprocal pairs are
\begin{align}
 J(1/3)&=\frac{\pi^2}{9}+\frac12\log^22-\frac12\log^23-\Li_2(1/3),
 \label{eq:J13}\\
 J(1/5)&=\frac{7\pi^2}{60}+\frac12\log^22-\frac14\log^25
                  -\log^2\phi-\frac12\Li_2(1/5).
 \label{eq:J15}
\end{align}
These are already present in the manuscript's $J(2/q)$ treatment
\cite{Repo}. For a direct derivation, use \eqref{eq:J-F} at $1/3$ and
$1/5$, the reciprocal-integer formula
\[
 F(1/m)-F(1)=-\frac{(m-1)(3m-2)}{24m}\pi^2-\frac12S(m),
\]
the three-term consequence
\[
 F(2/3)-F(1)=-\frac{5\pi^2}{36}+\Li_2(2/3)
                  +\frac12\log^2(3/2)-\frac12S(3)+\frac12S(2),
\]
and the displayed $F(2/5)$ formula. The elementary values of $S$ and
Euler reflection at $1/3$ and $1/5$ give
\eqref{eq:J13}--\eqref{eq:J15} without a branch change.

Equations \eqref{eq:J35}, \eqref{eq:J13}, \eqref{eq:J15} and
reciprocity evaluate every argument in \eqref{eq:six-extras}. Their
completeness as a list of formal exceptions is the new classification;
the first two classical evaluations are not newly asserted discoveries.

\begin{corollary}[Cancellation among the three exceptional ratios]
\label{cor:exceptional-relation}
\begin{equation}\label{eq:exceptional-relation}
 \boxed{\quad J(1/3)-J(1/5)+J(3/5)
      =\frac12\log^22+2\log^2\phi-\frac{2\pi^2}{45}.\quad}
\end{equation}
\end{corollary}
\begin{proof}
Substitute the three evaluations. Both rational dilogarithms, the
$\log^23$ term and the $\log^25$ term cancel. The rational
$\pi^2$ coefficient is $1/9-7/60-7/180=-2/45$.
\end{proof}
